\subsection{Paths}

DEFINITION 1. --- Let X be a topological space. A path in X is any continuous map c from I to X. The point \(c(0)\) is called the origin, the point \(c(1)\) the endpoint of the path c. A loop in X is a path in X whose origin is equal to the endpoint.

Let \(x\) and \(y\) be points of X. We say that a path \(c\) in X connects \(x\) to \(y\) if \(x\) is its origin and \(y\) its endpoint.

DEFINITION 2. --- Let X be a topological space. Paths c and d in X are said to be juxtaposable if one has \(c(1) = d(0)\). We then call the juxtaposed path of c and d the path \(c*d\) defined by the formula:

\[
(c * d) (t) = \left\{ \begin{array}{l l} c (2 t) & \text { for } 0 \leqslant t \leqslant 1 / 2, \\ d (2 t - 1) & \text { for } 1 / 2 \leqslant t \leqslant 1. \end{array} \right.\tag{1}
\]

Its origin is the origin of c, and its endpoint is that of d.

Let c be a path in X; one calls the path opposite to c, and denotes by \(\overline{c}\) the path defined by \(\overline{c}(t) = c(1 - t)\) for \(t \in I\).

If c and d are two juxtaposable paths in X, \(\overline{d}\) and \(\overline{c}\) are, and one has \(\overline{c*d} = \overline{d} * \overline{c}\) . For every path c in X, one has \(\overline{\overline{c}} = c\) .

Remarks. --- 1) Let X be a topological space and let P be a topological space reduced to a point. Identify P × I with I by the projection pr₂. The set C(P × I; X) of homotopies between maps (necessarily continuous) from P to X is then identified with the set C(I; X) of paths in X. To a homotopy connecting the constant map with image x to the constant map with image y there corresponds a path with origin x and endpoint y. The preceding identifications are compatible with the notions of juxtaposition and passage to the opposite (cf. III, p. 230).

\begin{enumerate}
\def\labelenumi{\arabic{enumi})}
\setcounter{enumi}{1}
\tightlist
\item
  Let X and Y be topological spaces. The canonical map
\end{enumerate}

\[
\mathcal {C} (\mathrm{X} \times \mathbf {I}; \mathrm{Y}) \rightarrow \mathcal {C} (\mathbf {I}; \mathcal {C} _ {c} (\mathrm{X}; \mathrm{Y}))
\]

associates to every homotopy \(\sigma\colon X \times I \to Y\) connecting two maps \(f\) and \(g\) from X to Y a path of origin \(f\), with endpoint \(g\), in the space \(\mathcal{C}_c(X;Y)\) . If \(X\) is a locally compact space, this canonical map is bijective (TG, X, p. 28, theorem 3).

Let X be a topological space. One calls the path space of X, and denotes by \(\Lambda(\mathrm{X})\) , the topological space \(\mathcal{C}_{c}(\mathbf{I};\mathrm{X})\) whose elements are the paths in X and whose topology is that of compact convergence (cf. TG, X, p. 27). If x is a point of X, one denotes by \(\Lambda_{x}(\mathrm{X})\) the subspace of \(\Lambda(\mathrm{X})\) formed by the paths with origin x. If y is a second point of X, one denotes by \(\Lambda_{x,y}(\mathrm{X})\) the subspace of \(\Lambda(\mathrm{X})\) formed by the paths with origin x and endpoint y.

PROPOSITION 1. --- Let X and Z be topological spaces. For a map \(\varphi\) from Z into the path space \(\mathcal{C}_c(\mathbf{I};\mathrm{X})\) to be continuous, it is necessary and sufficient that the map \((z,t)\mapsto \varphi (z)(t)\) be a continuous map from \(\mathbf{Z}\times \mathbf{I}\) into X. In particular, the map \((c,t)\mapsto c(t)\) of \(\mathcal{C}_c(\mathbf{I};\mathrm{X})\times \mathbf{I}\) into X is continuous.

The topological space I being locally compact, the proposition follows from TG, X, p. 28, th. 3.

One denotes by o and calls the origin map the map \(c \mapsto c(0)\) of \(\Lambda(\mathrm{X})\) into X; one denotes by e and calls the endpoint map the map \(c \mapsto c(1)\) of \(\Lambda(\mathrm{X})\) into X. The maps o and e are continuous (TG, X, p. 27, remark 1). The pairs of juxtaposable paths in X are the elements of the fibered product of the X-spaces \((\Lambda(\mathrm{X}), e)\) and \((\Lambda(\mathrm{X}), o)\) .

PROPOSITION 2. --- Let X be a topological space. The map which to a path c associates the opposite path \(\overline{c}\) is a homeomorphism of \(\Lambda(X)\) on itself. The juxtaposition of paths \((c,d)\mapsto c*d\) is a homeomorphism of the space \((\Lambda(X),e)\times_{X}(\Lambda(X),o)\) onto \(\Lambda(X)\) .

The map \(t \mapsto 1 - t\) is a homeomorphism of the space I onto itself; the first assertion follows from this.

Let C denote the fiber product \((\Lambda(\mathrm{X}),e)\times_{\mathrm{X}}(\Lambda(\mathrm{X}),o)\) . Denote by \(\gamma:C\to\Lambda(X)\) the map \((c,d)\mapsto c*d\) and \(\delta:C\times I\to X\) the map \(((c,d),t)\mapsto(c*d)(t)\) . The restrictions of the map \(\delta\) to \(C\times[0,\frac{1}{2}]\) and \(C\times[\frac{1}{2},1]\) are continuous by virtue of formula (1) and prop. 1. The map \(\delta\) is therefore continuous (TG, I, p. 19, prop. 4), as is the map \(\gamma\) (prop. 1). We have \(c(t)=(c*d)(\frac{t}{2})\) and \(d(t)=(c*d)(\frac{1+t}{2})\) , therefore \(\gamma\) is injective.

Let \(g\) be a path in \(\mathbf{X}\); for \(t \in \mathbf{I}\), set

\[
c _ {g} (t) = g \bigl (\frac {t}{2} \bigr) \quad \text{and} \quad d _ {g} (t) = g \bigl (\frac {1 + t}{2} \bigr).
\]

The maps \(c_{g}\) and \(d_{g}\) thus defined are juxtaposable paths in X and we have \(c_{g} * d_{g} = g\) . Moreover, the maps \(g \mapsto c_{g}\) and \(g \mapsto d_{g}\) are continuous maps of the space \(\Lambda(\mathrm{X})\) into itself (prop. 1). It follows that the map \(\gamma\) is a homeomorphism.

COROLLARY. --- Let X be a topological space and let x, y, z be points of X. The map \(c \mapsto \overline{c}\) from \(\Lambda_{x,y}(X)\) to \(\Lambda_{y,x}(X)\) and the map \((c,d) \mapsto c * d\) from \(\Lambda_{x,y}(X) \times \Lambda_{y,z}(X)\) to \(\Lambda_{x,z}(X)\) are continuous.

